\scanpage{19}
\mspage{9}
On a ainsi plusieurs points de vue essentiellement équivalents, pour décrire par voie profinie \unsure{\deleted{\unsure{une|d'une|des}}|\unsure{une|d'une|des}} extension \inserted{$K$} de type fini de $\mathbb Q$ :

1) Par le topos étale $B_K$, en tant que topos pro-galoisien sur $B_{\mathbb Q}$

2) Par le groupoïde fondamental $\Pi_K$ de \unsure{ce|en} topos, (le groupoïde de ses pts, ou de ses rev. universels) --- en tant que groupoïde au-dessus de $\Pi_{\mathbb Q}$

\deleted{2)} \deleted{\unsure{2}}3) Par le groupe extérieur $\unsure{\mathcal E|E}_K$, \unsure{au-dessus|--- disons|--- dessus} du groupe extérieur $\unsure{\mathcal E|E}_{\mathbb Q}\longleftarrow\Gamma_{\mathbb Q}$ (pro-gr)

4) En termes d'une clôture algébrique $\bar K/K$ (i.e. en décrivant le couple $(K,\bar K)$ plutôt que $K$), par un objet $\overline{\mathbb Q}\in\operatorname{Ob}\Pi_{\mathbb Q}$\unsure{,|)} et un hom de groupes profinis $\unsure{\mathcal E|E}\unsure{\longrightarrow|-}\Gamma_{\overline{\mathbb Q}/\mathbb Q}$.
\deleted{un ss-groupe ouvert $\Gamma$ de}

5) En termes d'une clôture alg. \inserted{\unsure{fixée|fixe}} $\overline{\mathbb Q}$ de $\mathbb Q$, \deleted{par le groupe}\inserted{et \unsure{si} $\Gamma=\Gamma_{\overline{\mathbb Q}/\mathbb Q}$}, \deleted{par} considérant les couples $(K,i)$ où $i:k\longrightarrow\overline{\mathbb Q}$ est un plongement de le \unsure{clôture|cloture} alg.\inserted{$k$} de $\mathbb Q$ dans $K$ dans $\overline{\mathbb Q}$ : \deleted{\unsure{comme|connexe|connexes}} par le groupe extérieur $\pi_K\unsure{\neq|=|\simeq}\pi_1(K)$, sur lequel un sous-groupe ouvert $\Gamma_K\subset\Gamma$ opère extérieurement.

\margin{6) En termes d'un $\overline{\mathbb Q}/\mathbb Q$ : par le \unsure{gr.-oïde|gr.-oïdes} $\Pi_{K\otimes_{\mathbb Q}\overline{\mathbb Q}}$ (pas néc. connexe) sur lequel $\Gamma_{\overline{\mathbb Q}/\mathbb Q}$ opère --- ou bien \unsure{plutôt|plutot}, par le topos quotient $B_{K\otimes_{\mathbb Q}\overline{\mathbb Q}}$\unsure{,|)} sous $\Gamma_{\overline{\mathbb Q}/\mathbb Q}$ opérant dessus\ldots}
\margin{18}
